\documentclass{amsart}
\usepackage{graphicx} % Required for inserting images
\usepackage{lipsum}
\usepackage{amsfonts, amsmath, amssymb, amsthm}
\title{Graphs and Matrices}
\usepackage{xcolor}
\usepackage[
colorlinks=true,
linkcolor=blue, 
citecolor=magenta,
urlcolor=red
]{hyperref}
\usepackage[sc]{mathpazo}
\usepackage[T1]{fontenc}
\author{Sayan Kar}
\date{March 2026}

\theoremstyle{plain}
\newtheorem{theorem}{Theorem}[section]
\newtheorem{corollary}{Corollary}[section]
\newtheorem{proposition}{Proposition}[section]
\begin{document}

\maketitle
\tableofcontents 

\begin{abstract}
    This is a very brief note on \textbf{Spectral Graph Theory}.
\end{abstract}

\section{Adjacency Matrices}

Given a graph, one can associate various matrices to encode its information. 
The \emph{adjacency matrix} $A$ of a graph $X = (V, E)$ is the matrix whose rows and columns 
are indexed by the vertices of $X$, where $A(x, y)$ equals the number of edges between 
$x$ and $y$. When necessary to indicate the dependence on $X$, we denote $A$ by $A(X)$.

A number $\lambda$ is called an \emph{eigenvalue} of $A$ if there exists a non-zero vector 
$u \in \mathbb{R}^n$ such that $ Au = \lambda u$. This means that $\lambda$ is a root of the \emph{characteristic polynomial} $\det(t I_n - A)$.

If $n$ denotes the number of vertices of $X$, then $A$ is an $n \times n$ real and 
symmetric matrix, and therefore by linear algebra, it has $n$ real eigenvalues, which we 
denote as follows: $\lambda_1 \ge \cdots \ge \lambda_n$.

The \emph{spectrum} of $X$ consists of its eigenvalues, including multiplicities.

\begin{theorem}
\label{clowalk}
Let $X = (V, E)$ be a graph without loops, with adjacency matrix $A$. 
For any $x, y \in V$ and any natural number $k$, the entry $A^k(x, y)$ 
equals the number of $(x, y)$-walks of length $k$.
\end{theorem}

\begin{proof}
Denote by $w_k(x, y)$ the number of $(x, y)$-walks of length $k$. 
We will show that $A^k(x, y) = w_k(x, y)$ for any $x, y \in V$ and any 
natural number $k$ using induction on $k$.

The base case $k = 1$ is clear from the definition of $A$. Let $k \ge 2$ and assume that $A^{k-1}(x, y) = w_{k-1}(x, y)$ for any 
$x, y \in V$. Let $x$ and $y$ be two vertices of $X$. A $(x, y)$-walk 
of length $k$ is composed of a $(x, z)$-walk of length $k-1$ followed 
by an edge $zy$. Hence,
\[
w_k(x, y) = \sum_{z : z \sim y} w_{k-1}(x, z).
\]

At the same time,
\[
A^k(x, y) = \sum_{z \in V} A^{k-1}(x, z) A(z, y) 
= \sum_{z : z \sim y} A^{k-1}(x, z).
\]

From the induction hypothesis, we know that 
$A^{k-1}(x, z) = w_{k-1}(x, z)$ for each $z \sim y$. 
Using the above two equations, we get that 
\[
A^k(x, y) = w_k(x, y).
\]
\end{proof}

Now the above thereom gives the next Corollary~\ref{clowalk} directly.

\begin{corollary}
Let $X = (V, E)$ be a graph with adjacency matrix $A$ whose eigenvalues are 
$\lambda_1 \ge \cdots \ge \lambda_n$. For any natural number $k$, the number 
of closed walks of length $k$ in $X$ equals $\sum_{j=1}^{n} \lambda_j^k$.
\end{corollary}

Now we shall see a very interesting upper bound on the eigenvalues of a simple graph through the following proposition.


\begin{proposition}
Let $X$ be a simple graph with $n$ vertices and $e$ edges. If $\lambda$ 
is an eigenvalue of the adjacency matrix $A$ of $X$, then
\[
|\lambda| \le \sqrt{2e\left(1-\frac{1}{n}\right)}.
\]
\end{proposition}

\begin{proof}
We use the facts that $\operatorname{tr}(A) = 0$ and $\operatorname{tr}(A^2) = 2e$. 
Since $A$ is real symmetric, all its eigenvalues are real. Let the eigenvalues of $A$ be $\lambda_1, \dots, \lambda_n$, and let $\lambda$ 
be one of them. The desired inequality is equivalent to
\[
\lambda^2 \le \frac{2e(n-1)}{n}
\;\;\Longleftrightarrow\;\;
\frac{n}{n-1}\lambda^2 \le \operatorname{tr}(A^2)
= \sum_{i=1}^n \lambda_i^2.
\]

Rewriting, this is equivalent to
\[
\sum_{i \ne k} \lambda_i^2 \ge \frac{1}{n-1}\lambda^2,
\]
where $\lambda = \lambda_k$. Now, since $\operatorname{tr}(A) = 0$, we have, $\sum_{i \ne k} \lambda_i = -\lambda$. Applying the Cauchy--Schwarz inequality,
\[
\sum_{i \ne k} \lambda_i^2 
\ge \frac{1}{n-1} \left( \sum_{i \ne k} \lambda_i \right)^2
= \frac{1}{n-1}\lambda^2.
\]

This proves the claim.
\end{proof}




\section{Bipartite Graphs}

A graph $X = (V, E)$ is called \emph{bipartite} if its vertex set $V$ can be partitioned 
into two disjoint sets $V_1$ and $V_2$ such that every edge of $X$ has one endpoint in 
$V_1$ and the other in $V_2$.

\begin{theorem}\label{bip_mul}
If $X$ is bipartite, and $\lambda > 0$ is an eigenvalue of the adjacency matrix $A$ 
of $X$ with multiplicity $m$, then $-\lambda$ is also an eigenvalue with multiplicity $m$.
\end{theorem}

\begin{proof}
Assume that $X$ is a bipartite graph whose partite sets have sizes $r$ and $s$, 
respectively. Label the vertices so that the adjacency matrix $A$ has the form
\[
A = \begin{pmatrix}
0_r & B \\
B^T & 0_s
\end{pmatrix},
\]
where $B$ is an $r \times s$ matrix.

Let $\lambda > 0$ be an eigenvalue of $A$ with eigenvector $w = \begin{pmatrix} u \\ v \end{pmatrix}$,
where $u \in \mathbb{R}^r$ and $v \in \mathbb{R}^s$, not both zero. Then
\[
Aw = \lambda w
\;\Longrightarrow\;
\begin{cases}
\lambda u = Bv, \\
\lambda v = B^T u.
\end{cases}
\]

Define, $w' = \begin{pmatrix} u \\ -v \end{pmatrix}.$. Then,
\[
Aw' = 
\begin{pmatrix}
0 & B \\
B^T & 0
\end{pmatrix}
\begin{pmatrix}
u \\ -v
\end{pmatrix}
=
\begin{pmatrix}
- Bv \\
B^T u
\end{pmatrix}
=
\begin{pmatrix}
-\lambda u \\
\lambda v
\end{pmatrix}
= (-\lambda) w'.
\]

Since $w \neq 0$, we also have $w' \neq 0$. Hence $-\lambda$ is an eigenvalue of $A$.

A collection of $m$ linearly independent eigenvectors corresponding to $\lambda$ 
gives rise to $m$ linearly independent eigenvectors corresponding to $-\lambda$. 
Thus the multiplicity of $-\lambda$ is at least that of $\lambda$. Repeating the 
argument with $-\lambda$ in place of $\lambda$ shows that the multiplicities are equal.
\end{proof}

\begin{theorem}[Spectral Characterization of Bipartite Graphs]
Let $X$ be a graph on $n$ vertices with adjacency matrix $A$ and eigenvalues 
$\lambda_1 \ge \cdots \ge \lambda_n$. The following statements are equivalent:
\begin{enumerate}
\item The graph $X$ is bipartite.
\item The spectrum of $X$ is symmetric with respect to zero.
\item $ t^{-(n\text{ mod }2)}\det\left( tI_n-A \right)\in \mathbb{R}[t^2].$
\item For any odd natural number $k$, $\sum_{j=1}^n \lambda_j^k = 0$.
\end{enumerate}
\end{theorem}

\begin{proof}
The implication (1) $\Rightarrow$ (2) was proved in Theorem~\ref{bip_mul}. To see that (2) $\Rightarrow$ (3), note that when $n$ is even, for any positive 
eigenvalue $\lambda$ of multiplicity $m$, we also have the eigenvalue $-\lambda$ 
with the same multiplicity $m$. Hence the factors
\[
(t - \lambda)^m (t + \lambda)^m = (t^2 - \lambda^2)^m,
\]
which is a polynomial in $t^2$. The eigenvalue $0$ has even multiplicity, and 
therefore $\det(tI_n - A)$ is a polynomial in $t^2$. When $n$ is odd, a similar argument applies, except that the multiplicity of $0$ 
is odd. Hence $t^{-1}\det(tI_n - A)$ is a polynomial in $t^2$. The implications (3) $\Rightarrow$ (2) and (2) $\Rightarrow$ (4) are trivially true. To see that (4) $\Rightarrow$ (1), note that condition (4) implies that the number 
of closed walks of odd length is zero. Therefore, $X$ has no odd cycles, and hence 
is bipartite. This completes the proof.
\end{proof}



\section{Upper Bound on Diameter}

The \emph{diameter} of a connected graph $X = (V, E)$ is defined as
\[
\operatorname{diam}(X) = \max_{x,y \in V} d(x,y).
\]
If $A$ is the adjacency matrix of $X$, the set of polynomials in $A$ with real 
coefficients is a vector space over $\mathbb{R}$. It is called the \emph{adjacency 
algebra} of $X$, and its dimension equals the number of distinct eigenvalues of $X$.

\begin{theorem}
If $X$ is a connected graph whose adjacency matrix has $r$ distinct eigenvalues, 
then $\operatorname{diam}(X) \le r - 1$.
\end{theorem}

\begin{proof}
Let $d = \operatorname{diam}(X)$. We show that $I_n, A, \dots, A^d$ are linearly 
independent matrices. Since the dimension of the adjacency algebra is $r$, this 
implies that $d + 1 \le r$, which proves the theorem.

To show that $I_n, A, \dots, A^d$ are linearly independent, we prove that for any 
$1 \le k \le d$, the matrix $A^k$ cannot be written as a linear combination of 
$I_n, A, \dots, A^{k-1}$.

Fix $1 \le k \le d$ and choose vertices $x, y \in V$ such that $d(x,y) = k$. 
Then there exists an $(x,y)$-path of length $k$, and there are no $(x,y)$-walks 
of length $\ell$ for any $\ell < k$. By Theorem~\ref{clowalk},
\[
A^k(x,y) > 0 \quad \text{and} \quad A^\ell(x,y) = 0 \text{ for all } \ell < k.
\]
Therefore, $A^k$ cannot be written as a linear combination of 
$I_n, A, \dots, A^{k-1}$. This completes the proof.
\end{proof}

This bound is actually sharp. Take $K_n$, the adjacency matrix $A = (A(x,y)) = J_n - I_n$, i.e., $A(x,y) = 1 - \delta(x,y)$. We know the characteristic polynomial of $J_n$ is $t^{n-1}(t-n)$, thus the characteristic polynomial of $A$ is $(t+1)^{n-1}(t - n+1)$, thus only two distinct eigenvalues of $A$ and diameter of $K_n$ is 1.


\section{Laplacian Matrices}

First consider $X$ as a directed graph. For each edge $e$, one of its endpoints becomes the \emph{tail} and the other endpoint is the \emph{head} of the directed edge.

Given such a direction on $X$, the \emph{directed incidence matrix} $N$ is defined 
as follows. Its rows are indexed by the vertex set $V$ and its columns by the edge 
set $E$. For $x \in V$ and $e \in E$,
\[
N(x,e) =
\begin{cases}
+1, & \text{if $x$ is the head of $e$}, \\
-1, & \text{if $x$ is the tail of $e$}, \\
0, & \text{otherwise}.
\end{cases}
\]

Also denote by $D$ the $V \times V$ diagonal matrix whose $(x,x)$-entry equals 
the degree of the vertex $x$. We call $D$ the \emph{diagonal degree matrix} of $X$. 

 The \emph{Laplacian matrix} $L$ of a graph $X$ is defined as $L = D - A$, where 
$D$ is the diagonal degree matrix of $X$ and $A$ is the adjacency matrix of $X$. 
This is another important matrix associated to a graph, and we summarize its basic 
properties below.

We have seen before that spectrum of adjacency matrix can determine whether a graph is bipartite or not but it can not guarantee connectedness of a graph. For example, take $C_4$ with an isolated vertex and $K_{1,4}$, both have the same spectrum but one is connected and the other is not.

\begin{theorem}
Let $X = (V, E)$ be a graph and let $N$ be a directed incidence matrix of $X$.
\begin{enumerate}
\item The Laplacian matrix $L$ equals $NN^T$.
\item The Laplacian matrix $L$ is real and symmetric, and its eigenvalues are nonnegative.
\item The smallest eigenvalue of $L$ is $0$, and its multiplicity equals the number of connected components of $X$.
\end{enumerate}
\end{theorem}

\begin{proof}
(1) Let $x, y \in V$. If $x = y$, then the $(x,x)$-entry of $L$ is $d(x)$. Also,
\[
(NN^T)(x,x) = \sum_{e \in E} N(x,e)N^T(e,x) 
= \sum_{e \in E} N^2(x,e) = d(x).
\]

If $x \ne y$ and $x \sim y$, then $L(x,y) = -1$. Let $f$ be the edge connecting 
$x$ and $y$. Regardless of its direction, $N(x,f)N(y,f) = -1$, while for any 
other edge $e \ne f$, $N(x,e)N(y,e) = 0$. Hence,
\[
(NN^T)(x,y) = \sum_{e \in E} N(x,e)N(y,e) = -1.
\]

If $x \ne y$ and $x \not\sim y$, then $L(x,y) = 0$. In this case, for any 
$e \in E$, $N(x,e)N(y,e) = 0$, and therefore
\[
(NN^T)(x,y) = 0.
\]
Thus, $L = NN^T$.

\medskip

(2) Since $D$ and $A$ are real and symmetric and $L = D - A$, it follows that 
$L$ is real and symmetric. Let $\mu$ be an eigenvalue of $L$ with eigenvector $u$. 
Then
\[
\mu u = Lu = NN^T u.
\]
Multiplying on the left by $u^T$, we obtain
\[
\mu\, u^T u = u^T NN^T u = (N^T u)^T (N^T u).
\]
Since $u \ne 0$, we have $u^T u = \sum_{x \in V} u_x^2 > 0$, and hence
\[
\mu = \frac{(N^T u)^T (N^T u)}{u^T u}
= \frac{\sum_{xy \in E} (u_x - u_y)^2}{\sum_{x \in V} u_x^2} \ge 0.
\]

\medskip

(3) Let $\mathbf{1}$ be the all-ones vector. Then each row of $L$ sums to zero, so 
$L\mathbf{1} = 0$. Hence $0$ is an eigenvalue of $L$. Assume $X$ is connected. Let $u$ be an eigenvector corresponding to $0$. Then
\[
0 = \frac{\sum_{xy \in E} (u_x - u_y)^2}{\sum_{x \in V} u_x^2},
\]
which implies that $u_x = u_y$ for every edge $xy \in E$. Since $X$ is connected, 
it follows that $u_x = u_y$ for all $x, y \in V$. Thus $u$ is a scalar multiple of 
$\mathbf{1}$, and the multiplicity of $0$ is $1$.

In general, $L$ can be written as a block diagonal matrix whose diagonal blocks are 
the Laplacian matrices of the connected components of $X$. The result follows.
\end{proof}

Now we have a sufficient condition for a graph to be connected. If the Laplacian Matrix of the graph has eigenvalue $0$ with multiplicity $1$, then the graph is connected.





\begin{thebibliography}{9}

\bibitem[RB]{bapat}
R. B. Bapat,
\textit{Graphs and Matrices},
Universitext, Springer, London, 2010.

\bibitem[RMC]{murty}
{M. Ram Murty, Sebastian M. Cioabă}, 
\textit{A First Course in Graph Theory and Combinatorics},
Springer, 2023.

\end{thebibliography}

\end{document}
